\section{\sys}
\label{sec:method}

\subsection{Preliminary}
\label{sec:prelim}

\paragraph{Agentic prompting.}
Let $x = [x_t]_{t=1}^{L}$ denote a committed trunk of length $L$, and let $\rho_0,\ldots,\rho_{k-1}$ denote $k$ residuals proposed from $x$.
Each candidate sequence is $z_i = x \circ \rho_i$.
After the step, a single spine $x \circ \rho_{\star}$ remains live.
$p_{\theta}(\cdot \mid z_i)$ is the next-token distribution of an LLM parameterized by $\theta$.
The model is expected to answer from at least one live residual, while abortable residuals should contribute neither prefill nor transfer.

State-of-the-art LLMs adopt the Transformer architecture~\citep{transformer}, where each token of $z_i$ is encoded by layer $\ell$ into a key and a value.
Let $H_i$ be the KV cache of $z_i$ after prefill.
For a continuation $y_i = [y_{i,t}]_{t=1}^{D}$ of budget $D$,
\begin{equation}
\label{eq:decode}
p_{\theta}(y_{i,t} \mid z_i, y_{i,<t}) = p(y_{i,t} \mid H_i, y_{i,<t}, \theta).
\end{equation}
$H_i$ is reused during decode with different $y_{i,t}$.
On a copy-on-write tree, $x$ is aliased in $O(1)$ and only $\rho_i$ allocates new pages; live occupancy after abort is the spine.

\paragraph{Costs of one fan-out.}
The interval that matters is from the moment the residuals are proposed until one of them is committed and the rest are aborted.
That interval contains three costs.
GPU prefill $W$ is the tokens the kernel actually runs.
The connector payload $X$ is the KV that a disaggregated prefill instance ships to a decode instance~\citep{distserve,splitwise}.
Decode-to-cut is the prefix that must be generated before a prune score exists.
On the systems in Figure~\ref{fig:problem}, all three are paid before the abort.

\paragraph{Decode-time pruning.}
Early-stopping self-consistency (ESC), speculative rejection, and dynamic parallel tree search (DPTS) score tokens that are already resident~\citep{esc,specrej,dpts}.
As illustrated in Figure~\ref{fig:problem}(b), we define $y_i$ as a continuation of $z_i$ under a host cut: ESC decodes to $D$ while an agreement window is open, speculative rejection scores at a partial-reward horizon, and DPTS generates a mini-step.
A residual is wasted if $H_i$ is computed and $y_i$ is later aborted.
An LLM serving stack is inefficient on the fan-out if abortable residuals are preferred for encoding over skipping, i.e., if they enter the prefill batch.

\subsection{Refusing Abortable Residuals}
\label{sec:refuse}

We refuse by steering the stack to interpret abortable residuals as skippable text rather than as thoughts to encode, thereby reducing $W$ and $X$ in response generation.
To achieve this, we propose \sys in two stages, Prefill Admission and Decode Continuation, illustrated in Figure~\ref{fig:arch}.

\begin{figure*}[t]
\centering
\resizebox{\textwidth}{!}{\begin{tikzpicture}[
  x=1cm, y=1cm,
  font=\sffamily\scriptsize,
  >=Latex,
  line cap=round, line join=round,
  panel/.style={draw=black!20, line width=0.55pt, rounded corners=2.5pt, fill=white},
  head/.style={font=\scriptsize\bfseries\sffamily, text=white, anchor=west},
  box/.style={draw, rounded corners=1.6pt, inner sep=2.2pt, font=\tiny\sffamily,
              line width=0.45pt, align=center, minimum height=0.42cm},
  llm/.style={box, draw=fsblue!70, fill=fslight, minimum width=2.15cm, minimum height=0.72cm},
  kv/.style={box, draw=fsgray!55, fill=black!4, minimum width=1.55cm},
  arr/.style={-{Latex[length=1.7mm,width=1.3mm]}, line width=0.62pt, draw=fsblue!75},
]
\path (0,0) rectangle (16.40,8.05);

% ===== query / fan-out =====
\draw[panel, fill=black!3] (0.00,7.28) rectangle (16.40,8.05);
\node[box, draw=fsblue!60, fill=fslight, anchor=west] at (0.18,7.66) {Query};
\draw[arr] (1.55,7.66) -- (2.05,7.66);
\fill[fsblue!75] (2.10,7.44) rectangle (5.35,7.88);
\node[font=\tiny\sffamily, text=white] at (3.72,7.66) {trunk $x$ of length $L$};
\draw[arr] (5.40,7.66) -- (5.90,7.66);
\node[box, fill=fscommit, draw=fsgreen!65!black] at (6.85,7.66) {$\rho_0$ viable};
\node[box, fill=fsmiss, draw=fsred!65] at (8.55,7.66) {loop};
\node[box, fill=fsspec, draw=fsorange!80!black] at (10.25,7.66) {illegal};
\node[box, fill=fscommit, draw=fsgreen!65!black] at (12.05,7.66) {viable};
\node[font=\tiny\sffamily, text=fsgray, anchor=west] at (13.10,7.66) {$\rho_0$ always admitted};

% ===== three stages =====
\draw[panel, fill=fsspec!35] (0.00,1.92) rectangle (5.15,7.12);
\draw[panel, fill=fslight] (5.45,1.92) rectangle (10.75,7.12);
\draw[panel, fill=fscommit] (11.05,1.92) rectangle (16.40,7.12);
\fill[fsorange!88!black] (0.00,6.68) rectangle (5.15,7.12);
\fill[fsblue] (5.45,6.68) rectangle (10.75,7.12);
\fill[fsgreen!75!black] (11.05,6.68) rectangle (16.40,7.12);
\node[head] at (0.14,6.90) {Stage I: Admit};
\node[head] at (5.59,6.90) {Stage II: Prefill};
\node[head] at (11.19,6.90) {Stage III: Probe and Resume};

% --- lane y: viable 5.85, loop 4.35, illegal 2.85 ---
% Stage I scorer
\node[box, draw=fsorange!75, fill=fsspec, minimum width=4.55cm, anchor=north]
  at (2.575,6.55) {$s^{\mathrm{pre}}(\rho_i)$ on residual text \quad vs.\quad $\tau_{\mathrm{pre}}$};

% viable lane
\node[box, fill=fscommit, draw=fsgreen!65!black, minimum width=2.05cm] at (1.55,5.85) {viable};
\draw[arr, draw=fsgreen!55!black] (2.65,5.85) -- (3.15,5.85);
\node[box, fill=white, draw=fsgreen!65!black, minimum width=1.35cm] at (3.95,5.85) {keep};

% loop lane
\node[box, fill=fsmiss, draw=fsred!65, minimum width=2.05cm] at (1.55,4.35) {loop};
\draw[arr, draw=fsred!70] (2.65,4.35) -- (3.15,4.35);
\node[box, fill=white, draw=fsred!65, minimum width=1.35cm] at (3.95,4.35) {drop};

% illegal lane
\node[box, fill=fsspec, draw=fsorange!80!black, minimum width=2.05cm] at (1.55,2.85) {illegal};
\draw[arr, draw=fsorange!75] (2.65,2.85) -- (3.15,2.85);
\node[box, fill=white, draw=fsorange!80!black, minimum width=1.35cm] at (3.95,2.85) {keep};

\node[font=\tiny\sffamily, text=fsgray, align=left, anchor=west] at (0.22,2.22)
  {$\mathrm{keep}_i=[i{=}0\lor s^{\mathrm{pre}}{\ge}\tau_{\mathrm{pre}}]$. No kernel.};

% arrows I -> II
\draw[arr] (5.15,5.85) -- (5.45,5.85);
\draw[arr, draw=fsred!55, densely dashed] (5.15,4.35) -- (5.45,4.35);
\draw[arr] (5.15,2.85) -- (5.45,2.85);

% Stage II
\node[llm] at (6.85,5.85) {LLM (frozen)};
\draw[arr] (8.00,5.85) -- (8.45,5.85);
\node[kv] at (9.45,5.85) {KV $H_{\star}$};

\node[draw=fsred!45, densely dashed, rounded corners=1.6pt, minimum width=4.70cm,
      minimum height=0.72cm] at (8.10,4.35) {};
\node[font=\tiny\sffamily, text=fsred] at (8.10,4.35) {withheld: $W{=}X{=}0$};

\node[llm] at (6.85,2.85) {LLM (frozen)};
\draw[arr] (8.00,2.85) -- (8.45,2.85);
\node[kv] at (9.45,2.85) {KV $H_{i}$};

\node[font=\tiny\sffamily, text=fsgray, align=left, anchor=west] at (5.62,2.22)
  {Kernel and transfer only on $\mathrm{Keep}$.};

% arrows II -> III
\draw[arr, draw=fsgreen!55!black] (10.75,5.85) -- (11.05,5.85);
\draw[arr, draw=fsorange!70] (10.75,2.85) -- (11.05,2.85);

% Stage III
\node[box, draw=fsgreen!60, fill=fscommit, minimum width=2.35cm] at (12.40,5.85) {decode to $D$};
\draw[arr, draw=fsgreen!55!black] (13.62,5.85) -- (14.12,5.85);
\node[box, draw=fsgreen!65!black, fill=white, minimum width=1.55cm] at (15.10,5.85) {resume};

\node[font=\tiny\sffamily, text=fsgray] at (13.72,4.35) {(no continuation)};

\node[box, draw=fsorange!75, fill=fsspec, minimum width=2.35cm] at (12.40,2.85) {probe $t_0$};
\draw[arr, draw=fsorange!75] (13.62,2.85) -- (14.12,2.85);
\node[box, draw=fsred!65, fill=white, minimum width=1.55cm] at (15.10,2.85) {stop};

\node[font=\tiny\sffamily, text=fsgray, align=left, anchor=west] at (11.22,2.22)
  {Extend iff $s^{\mathrm{dec}}(y_{1:t_0})\ge\tau_{\mathrm{dec}}$.};

% ===== number lines =====
\draw[panel, fill=black!2] (0.00,0.00) rectangle (16.40,1.76);
\node[font=\tiny\sffamily\bfseries, anchor=west] at (0.16,1.48)
  {The two thresholds are not interchangeable};

\draw[fsorange!80!black, line width=0.75pt] (0.95,0.62) -- (7.85,0.62);
\node[font=\tiny\sffamily, anchor=east] at (0.88,0.62) {prefill};
\foreach \x in {0.95, 1.85, 3.55, 4.85, 7.85} {
  \draw[fsorange!80!black] (\x,0.52) -- (\x,0.72);
}
\node[font=\tiny\sffamily, text=fsred] at (1.85,0.32) {loop};
\node[font=\tiny\sffamily, text=fsorange!80!black] at (3.55,1.00) {$\tau_{\mathrm{pre}}$};
\node[font=\tiny\sffamily, text=fsorange!80!black] at (4.95,0.32) {illegal};
\node[font=\tiny\sffamily, text=fsgreen!60!black] at (7.45,1.00) {viable};
\draw[fsorange!80!black, densely dashed] (3.55,0.28) -- (3.55,1.18);

\draw[fsblue, line width=0.75pt] (9.35,0.62) -- (16.05,0.62);
\node[font=\tiny\sffamily, anchor=east] at (9.28,0.62) {decode};
\foreach \x in {9.35, 10.10, 11.55, 13.85, 15.85} {
  \draw[fsblue] (\x,0.52) -- (\x,0.72);
}
\node[font=\tiny\sffamily, text=fsred] at (10.10,0.32) {loop};
\node[font=\tiny\sffamily, text=fsorange!80!black] at (11.55,0.32) {illegal};
\node[font=\tiny\sffamily, text=fsblue] at (13.85,1.00) {$\tau_{\mathrm{dec}}$};
\node[font=\tiny\sffamily, text=fsgreen!60!black] at (15.70,0.32) {answer};
\draw[fsblue, densely dashed] (13.85,0.28) -- (13.85,1.18);
\end{tikzpicture}
}
\caption{Illustration of the workflow of \sys.
Given a fan-out with live strategies, a loop, and illegal text, prefill scoring is computed from residual text before any kernel.
A residual is pruned by withholding it from the context KV cache (orange).
Activations after admission are generated only on the admitted cache.
Continuation uses a separate score on a generated prefix of length $t_0$ (green).
The number lines show why the two bars are not interchangeable: illegal text sits above $\tau_{\mathrm{pre}}$ and a still-illegal prefix sits below $\tau_{\mathrm{dec}}$.}
\label{fig:arch}
\end{figure*}

\subsubsection{Prefill Admission}
\label{sec:prefill-admit}

As mentioned in Section~\ref{sec:intro}, we aim at identifying residuals whose text biases the stack toward encoding the same span as a thought rather than as skippable text.
These residuals should not interfere with the model's ability to follow live strategies, i.e., by generating accurate answers that leverage the trunk as supporting context.
The identified residuals are excluded from the KV cache through the following steps.

\paragraph{Scoring.}
Suppose there exists a prefill score $s^{\mathrm{pre}}: \rho \mapsto [0,1]$ that captures the stack's preference over which residuals to encode, i.e., by measuring how much a residual looks like a live strategy rather than a loop.
Residuals scoring below $\tau_{\mathrm{pre}}$ can be viewed as steering the kernel toward wasted prefill.
We score only the residual span $\rho_i$, not the trunk $x$, because abortable text is proposed in the residual.
Details of $s^{\mathrm{pre}}$ are in Section~\ref{sec:scores}.

\paragraph{Forced keep.}
Index $0$ is the residual the search will keep if nothing else survives.
Pruning it would degrade the spine even when every sibling is abortable.
We therefore admit
\begin{equation}
\label{eq:keep}
\mathrm{keep}_i = [\, i = 0 \ \lor\  s^{\mathrm{pre}}(\rho_i) \ge \tau_{\mathrm{pre}} \,].
\end{equation}
Only residuals with $\mathrm{keep}_i$ are queued for the kernel.
$W$ is the measured prefill of that set.
With prefix caching on, a shared trunk is not the dominant term; the gate's effect is that a dropped residual contributes nothing.
The connector inserts KV only for the admitted set, so a dropped residual contributes $X_i = 0$ as well.

\paragraph{Thresholding.}
Pruning residuals with small $s^{\mathrm{pre}}$ should alter the stack's tendency to encode abortable text.
However, this would degrade live answers if $s^{\mathrm{pre}}$ were reused as a decode bar, because that score does not observe a generated prefix.
To maintain quality after pruning, we regularize by a second bar $\tau_{\mathrm{dec}} > \tau_{\mathrm{pre}}$ that is never applied to residual text in the default policy.
Raising the prefill bar to $\tau_{\mathrm{dec}}$ is a faster operating point that we measure, not the default.

\subsubsection{Decode Continuation}
\label{sec:decode-cont}

In experiments, the prefill mask is computed from residual text with no model call, then applied to every fan-out at test time.
Admitted residuals that the decode stage later stops have already paid their prefill; they do not pay the rest of $D$.

Decode length is chosen after admission.
A repeated loop is not probed: it never starts.
Index $0$ and any residual whose planner prior is already at least $\tau_{\mathrm{dec}}$ decode to the budget $D$.
Every other admitted residual decodes $t_0$ tokens, and is extended to $D$ only when a decode score $s^{\mathrm{dec}}$ on the generated prefix clears $\tau_{\mathrm{dec}}$:
\begin{equation}
\label{eq:extend}
\mathrm{extend}_i = [\, s^{\mathrm{dec}}(y_{i,1:t_0}) \ge \tau_{\mathrm{dec}} \,].
\end{equation}
Algorithm~\ref{alg:grow} is this rule.
Section~\ref{sec:hosts} is the special case in which a decode-time host, rather than $t_0$, supplies the later cut.

\begin{algorithm}[t]
\caption{Probe, then extend. Loops skip prefill. Index $0$ decodes to $D$.}
\label{alg:grow}
\begin{algorithmic}[1]
\Require residuals $\rho_{0:k-1}$, budget $D$, probe $t_0$, bars $(\tau_{\mathrm{pre}},\tau_{\mathrm{dec}})$
\For{$i \gets 0$ \textbf{to} $k-1$}
  \If{$i \neq 0$ \textbf{and} $\rho_i$ is a repeated loop}
    \State skip \Comment{no kernel, $X_i{=}0$}
  \ElsIf{$i = 0$ \textbf{or} planner prior of $\rho_i \ge \tau_{\mathrm{dec}}$}
    \State prefill $\rho_i$; decode to $D$
  \Else
    \State prefill $\rho_i$; decode $t_0$ tokens
    \If{$s^{\mathrm{dec}}(y_{i,1:t_0}) \ge \tau_{\mathrm{dec}}$}
      \State extend from $t_0$ to $D$
    \EndIf
  \EndIf
\EndFor
\end{algorithmic}
\end{algorithm}

Applying this mask over the residual ensures that the stack's implicit criteria align with the live--abortable separation already present in the text.
Varying the prefill bar, as in Section~\ref{sec:eval}, shows a trade-off between robustness of any-finished accuracy and quality of following only the winner.

\subsection{The Two Scores}
\label{sec:scores}

As described in Section~\ref{sec:refuse}, admission is guided by $s^{\mathrm{pre}}$ and continuation by $s^{\mathrm{dec}}$.
Both are deterministic functions of a string.
They do not call the model.
They do not share weights.

The prefill scorer looks only at $\rho_i$.
A repeated span is a loop and falls below $\tau_{\mathrm{pre}}$.
A residual that matches an illegal-math pattern, such as a run of stars or an undefined token, lands strictly between $\tau_{\mathrm{pre}}$ and $\tau_{\mathrm{dec}}$.
Any other non-empty residual receives
\begin{equation}
\label{eq:prefill-score}
s^{\mathrm{pre}}(\rho_i) = 0.55\,\min(1, 2\alpha_i) + 0.45\,\min(1, 8\eta_i),
\end{equation}
where $\alpha_i$ is the alphanumeric fraction and $\eta_i$ is the character-level uniqueness.
Live strategy text clears both bars.

The decode scorer looks only at a generated prefix $y_i$.
A repeated prefix stays below $\tau_{\mathrm{dec}}$.
An answer mark (\texttt{\#\#\#\#} or \texttt{\textbackslash boxed}) clears $\tau_{\mathrm{dec}}$.
Illegal text on a prefix also stays below $\tau_{\mathrm{dec}}$, which is why it cannot share $\tau_{\mathrm{pre}}$.
Otherwise
\begin{equation}
\label{eq:decode-score}
\begin{aligned}
s^{\mathrm{dec}}(y_i) = {} & 0.50\,\min(1, 2\alpha) + 0.30\,\min(1, 8\eta) \\
& + 0.20\,\min(1, 20\delta + 15\omega),
\end{aligned}
\end{equation}
with $\delta$ the digit fraction and $\omega$ the fraction of arithmetic symbols.
A static planner prior used by the host-side ranker is a third mix again ($0.35$ alphanumeric, $0.65$ uniqueness) and is never reused as the prefill score.

The operating point is the ordered pair $(\tau_{\mathrm{pre}}, \tau_{\mathrm{dec}})$ with $\tau_{\mathrm{pre}} < \tau_{\mathrm{dec}}$ (Equation~\ref{eq:pair} in the appendix gives the run).
On the prefill axis the loop is below $\tau_{\mathrm{pre}}$ and illegal text is above it, so the gate drops loops and keeps illegal text.
On the decode axis illegal text is below $\tau_{\mathrm{dec}}$ and an answer mark is above it, so a probe stops on illegal text and continues once an answer mark appears.
Section~\ref{sec:eval} measures what happens when a single bar is used for both decisions.

\paragraph{The subset that must survive.}
In Section~\ref{sec:prefill-admit}, we regularize admission by forcing index $0$ and by not using $\tau_{\mathrm{dec}}$ as the default prefill bar.
We want to avoid pruning residuals whose prefill score is that of live strategy text, so that the pruned stack can still generate answers that address the user query by leveraging the trunk.
Index $0$ is among the live set in our evaluation mix, so the forced keep does not add a further slot beyond the live half.
A batch-limited rule that ignores the score and keeps a single residual is a different policy; Section~\ref{sec:eval} shows that it discards correct branches.

\begin{proposition}[Host-invariant prefill]
\label{prop:invariance}
Let $\mathrm{Keep} = \{ i : \mathrm{keep}_i \}$ be the admitted set of \eqref{eq:keep}.
For any decode-time host that scores only generated tokens, the prefill work $W$ equals the measured prefill of $\mathrm{Keep}$ and does not depend on the host.
\end{proposition}
The proof is in Appendix~\ref{app:invariance}.
Proposition~\ref{prop:invariance} is the reason three otherwise different hosts share one prefill column in Figure~\ref{fig:problem}(b) and in Table~\ref{tab:hosts}.

\subsection{The Score Gap}
\label{sec:gap}

Prefill scoring with \eqref{eq:prefill-score} relies on the residual string alone.
Not every abortable residual is a loop.
Specifically, we find that the same illegal span can be interpreted as a thought to encode or as a prefix to stop, depending on whether the score is read before or after generation.
We refer to this phenomenon as the \emph{score gap}.

To illustrate, illegal text scores above $\tau_{\mathrm{pre}}$ on the residual and below $\tau_{\mathrm{dec}}$ on a generated prefix that is still illegal.
A loop is below both bars.
A live residual clears $\tau_{\mathrm{pre}}$; a prefix that has emitted an answer mark clears $\tau_{\mathrm{dec}}$.
Figure~\ref{fig:arch} plots these four points on two number lines.
The gap is why a single bar is the wrong cut: applying $\tau_{\mathrm{dec}}$ as a prefill bar drops illegal residuals that can still solve once a prefix is visible, and applying $\tau_{\mathrm{pre}}$ as a decode bar would extend an illegal prefix that the decode rule is designed to stop.

Motivated by the gap, we perform continuation using only a short prefix of length $t_0$, which has been enough to capture the difference between an illegal prefix and an answer-bearing one.
Compared with decoding every admitted residual to $D$, the probe improves fidelity, since spending the full budget on an illegal prefix introduces additional decode work if a few tokens already suffice to govern the decision.
In experiments, $t_0{=}128$ is the setting that returns MATH-500 any-finished accuracy at $k{=}16$ to the full fan-out.
We do not assume testing prompts contain an answer mark at the first token.

\subsection{In Front of a Decode-Time Host}
\label{sec:hosts}

ESC, speculative rejection, and DPTS remain the decode policy.
\sys only changes the set of residuals that reach them.
\begin{itemize}
  \item ESC decodes an admitted residual out to $D$ while its agreement window is open.
    On this mix the window stays open, so an admitted residual runs to $D{=}512$ and a residual refused by \eqref{eq:keep} is never decoded.
  \item Speculative rejection scores at horizon $\tau{=}256$ and drops the lower fraction $\alpha{=}0.5$ of the admitted set.
    The memory-limit trigger from that paper does not fire at these widths, so $256$ is the cut.
  \item DPTS generates a mini-step of $100$ tokens and early-stops a child whose confidence is below the decode bar.
\end{itemize}
The prefill bar is not retuned per host.
$(\tau_{\mathrm{pre}}, \tau_{\mathrm{dec}})$ is the same pair on all three.
When the prefill bar is raised to $\tau_{\mathrm{dec}}$, the admitted set is the live half.
Those residuals are not the ones the host would have cut, the decode-side drop count falls to zero, and the three hosts converge to the same decode length.
The host remains complementary to \sys, since our approach does not modify the host's scoring rule.
